% ----------------------
\subsection{Section 27 (circa August 1830)}
% ----------------------
	\label{DC27}
	
	% -----
	\subsubsection{Changes to this section and versions of it}
	% -----
	
		The current version of this section is 18 verses long. The original version, found in the JSP and originating in RB1, is about a third of that length; it has verses 1--4, part of 5, 14, part of 15, and part of 18. 
		
		BC is basically identical to JSP/RB1.
		
		The additions which make the section into what we know now first appear in the 1835 DC.
		
		Note that this revelation also appeared in EMS, March 1833. That version is basically identical to JSP/RB1 as well, and so it's basically identical to the BC version, but the paragraph breaks are slightly different between the BC and EMS versions.
		
		Parts of this revelation also appeared in early versions of the PGP. See HDDC 395--397 for more.
		
		There is an early version of this revelation in the handwriting of Edward Partridge (mentioned at HDDC 395), and it was also published in the \textit{Painesville Telegraph} at one point. JSP has this to say: ``Two other versions, one copied by Edward Partridge (with minor gaps in the text of the manuscript) and the other published in the \textit{Painesville (OH) Telegraph}, are contemporaneous to the featured text below.1 No significant variant reading distinguishes any of the three texts as being clearly produced earlier than the other two. But because one of the other versions is missing text and the other includes added punctuation, the text below [from RB1] is likely the best representation of the earliest text'' (D1:164).
	
	% -----
	\subsubsection{Historical Background}
	% -----
		\label{DC27-HB}
		
		% --
		\paragraph{JSR}
		% --
		
			JSR has an extended commentary on the background of this section, the changes, and the title ``Elias,'' among other things. See that commentary \hyperref[EXSS-DC27-JSR]{here}.
		
		% --
		\paragraph{JSP}
		% --
		
			``This revelation announced guidelines for what members of the Church of Christ should use in the sacrament of the Lord's Supper to represent the blood of Christ. JS's history explained that he dictated the revelation during a visit from Newel and Sally Knight to his home in Harmony in August 1830: `As neither his wife nor mine had been as yet confirmed, it was proposed that we should confirm them, and partake together of the sacrament, before he and his wife should leave us.--- In order to prepare for this; I set out to go to procure some wine for the occasion, but had gone only a short distance when I was met by a heavenly messenger, and received the following revelation.'%
				\footnote{%
					% \label{}%
					This is from volume A-1 of the history.
					}
			The revelation told JS to use only new wine made by members of the church.
			
			``Both the dating and the text of this revelation present challenges. The earliest extant copy, in Revelation Book 1 (the version featured here), dated it broadly to the year 1830, and John Whitmer placed it between revelations dated July and September 1830. Despite Whitmer's placement, however, the versions published in 1833 in The Evening and the Morning Star and the Book of Commandments specified 4 September 1830 as the date. The versions published in the 1835 Evening and Morning Star reprint and in the 1835 edition of the Doctrine and Covenants included two additional paragraphs and dated the entire revelation to September 1830, removing the precise date mentioned in the earlier publications. In his later history, however, JS said that the first paragraph of the 1835 text `was written at this time [early August 1830], and the remainder in the September following.'%
				\footnote{%
					% \label{}%
					Also from A-1 of the history.
					}
			JS, then, affirmed that the first part should be dated August 1830 and the remainder September 1830, in which case the date of 4 September 1830 found in the earlier printed versions may have reflected the date of dictation of the second portion.
			
			``Early manuscripts, including the copy featured here, contain only the first portion of the combined revelation as published in 1835. Although an earlier manuscript may have existed for the September portion, the earliest extant text for the expanded version of the revelation is the 1835 edition of the Doctrine and Covenants. For the annotated treatment of the expanded version, see the volume of the Documents series that covers 1835''%
				\footnote{%
					% \label{}%
					See JSP D4:408--412. See \hyperref[EMS-1835-JS-CIR-AUG-1835]{here} for the JSP text, which comes from the 1835 DC, and below for more background.
					} 
			(D1:165).
			
			``In August and September 1830, JS dictated a revelation providing guidelines about the food and drink that were to be used in the administration of the sacrament. The revelation was recorded in Revelation Book 1, likely in 1831, and published, with minor edits, in the March 1833 issue of The Evening and the Morning Star and the 1833 Book of Commandments. The revelation was published again in the 1835 Doctrine and Covenants. The 1835 version, featured here, contains information not present in the earlier versions; in fact, less than one third of the 1835 text appears in the 1831 or 1833 texts. The additional information contains considerable detail about Jesus Christ one day partaking of sacramental wine with JS and various prophets and apostles from the Bible and the Book of Mormon. Also included is material emphasizing the transmission of priesthood keys---or the authority to govern and lead the church---to JS by biblical prophets, apostles, and patriarchs. In particular, the revelation outlines specific keys held by Elias and Elijah and references a visit by Peter, James, and John to ordain JS and Oliver Cowdery as apostles. This is the first clear documentary reference to angelic visitations by Peter, James, and John and the first clear reference to their conferral of priesthood keys.%
				\footnote{%
					% \label{}%
					JSP note: ``The date of Peter, James, and John's visit is unknown, but sources indicate it occurred sometime after John the Baptist's May 1829 visit to confer the lesser priesthood on JS and Oliver Cowdery. A later JS history indicates that John the Baptist acted under the direction of Peter, James, and John and promised they would later provide a higher, or additional, authority. In the preface to his 1832 history, JS indicated his intention to record his `reception' of both `the holy Priesthood' and `the high Priesthood,' neither of which the unfinished 1832 history addressed. Cowdery also noted receiving the Melchizedek, or `high and holy,' Priesthood in Patriarchal Blessing Book 1, but the featured text predates that entry by a month'' (D4:408--409). For the reference to Cowdery and the Patriarchal Blessing Book, see EPB 3, which records Cowdery's introductory comments to some early blessings; these comments were apparently written in September 1835: ``But before baptism, our souls were drawn out in mighty prayer to know how we might obtain the blessings of baptism and of the Holy Spirit, according to the order of God, and we diligently sought for the right of the fathers and the authority of the holy priesthood, and the power to admin[ister] in the same: for we desired to be followers of righteousness and the possessors of greater knowledge, even the knowledge of the mysteries of the kingdom of God. Therefore, we repaired to the woods, even as our father Joseph said we should, that is to the bush, and called upon the name of the Lord, and he answered us out of the heavens, and while we were in the heavenly vision the angel came down and bestowed upon us this priesthood; and then, as I have said, we repaired to the water and were baptized. After this we received the high and holy priesthood: but an account of this will be given elsewhere, or in another place.'' Note how some of Cowdery's language here is similar to the Book of Abraham, which they were working on at this same time; the bits about ``sought for the right of the fathers and the authority of the holy priesthood,'' and other parts. See, for example, the \hyperref[EMS-1835-JS-CIR-JUL-CIR-NOV-1835]{Book of Abraham manuscript} from the second half of 1835.
					}
			
			``The historical record is silent as to when or how the additional information in the 1835 version was originally recorded, though JS was undoubtedly involved. The absence of the content in the extant 1831 or 1833 versions suggests the extra information may have originated sometime thereafter. In addition, the 1835 version of the revelation identifies the archangel Michael as Adam and equates Adam with the ancient of days referred to in the Book of Daniel. Not until at least 1833 did church members appear to have identified Adam with Michael and the ancient of days.%
				\footnote{%
					% \label{}%
					JSP note: ``In January 1834, Cowdery wrote a letter to John Whitmer in which he said that since coming to Kirtland, Ohio, in summer 1833, he had learned `that the Angel Michael is no less than our father Adam.' Sometime around spring 1835, JS prepared an Instruction on Priesthood that also contained information about Adam being `Michael, the Prince, the Archangel.' Similarly, in June 1835, William W. Phelps wrote in a letter to Cowdery that he understood that `Michael, the prince' was `our great father Adam,' something that he considered to be `new light.' Around September 1835, Cowdery recorded a December 1833 blessing that JS gave to Joseph Smith Sr. and also referred to Adam as `Michael, the Prince, the Arch angel,' and `the Ancient of Days,' language that was not used in the original version of the blessing as recorded in JS's journal in 1833.''
					}
			However, it is possible that JS dictated all of this information in 1830. JS's history indicates that the `first paragraph' of the revelation, which represents most of the original revelation, was written down immediately in early August 1830 and the `remainder in the September following.' But if all of what appears in the 1835 version was originally dictated in 1830, it is unclear why the additional material was not included in the extant 1831 and 1833 versions.
			
			``As published in the 1835 edition of the Doctrine and Covenants, the revelation contains much information about the conferral of keys, or authority, upon JS. Although some 1829 and 1830 revelations spoke of JS holding keys, they referred specifically to the keys to bring forth the Book of Mormon and other ancient records. In 1831 and 1832, revelations referred frequently to JS and others holding keys that allowed them to administer the church. For example, revelations indicated that JS held `the keys of the mysteries of the kingdom' and `the keys of the kingdom of God,' similar to `the keys of the kingdom of heaven' held by Peter in the New Testament. This 1835 version of the revelation expands on these concepts, specifying what keys had been given to JS and others, who had provided those keys, and, by inference, why the keys were necessary.
			
			``The additions published in the 1835 Doctrine and Covenants consist of three intertextual insertions to the second paragraph of the original text'' (D4:408--410).
			
		% --
		\paragraph{HDDC}
		% --
		
			``Section 27 is either a composite of two revelations, or one revelation written in two parts. The uncertainty concerning its origin can be traced to two contemporary accounts. The first is given in the \uline{History of the Church}, wherein one introduction has been written for both parts as follows:
			
				\begin{quote}
					Early in the month of August Newel Knight and his wife paid us a visit at my place in Harmony, Pennsylvania; and as neither his wife nor mine had been yet confirmed, it was proposed that we should confirm them, and partake together of the Sacrament, before he and his wife should leave us. In order to prepare for this I set out to procure some wine for the occasion, but had gone only a short distance when I was met by a heavenly messenger, and received the following revelation, the first four paragraphs of which were written at this time, and the remainder in the September following.
				\end{quote}
				
			``Newel Knight was aware that this revelation was given during his visit, and wrote the following account in his diary:
			
				\begin{quote}
					In the beginning of August I, in company with my wife, went to make a visit to Brother Joseph Smith, Jn., who then resided at Harmony Penn. We found him and his wife well, and in good spirits. We had a happy meeting. It truly gave me joy to again behold his face. As neither Emma, the wife of Joseph Smith, nor my wife had been confirmed, we concluded to attend to that holy ordinance at this time, and also to partake of the sacrament, before we should leave for home. In order to prepare for this, Brother Joseph set out to procure some wine for the occasion, but he had gone only a short distance, when he was met by a heavenly messenger, and received the first four verses of the revelation given on page 138, of the Doctrine and Covenants (new Edition), the remainder being given in the September following at, Fayette, New York.%
						\footnote{%
							% \label{}%
							HDDC's reference here is: ``Journal History of the Church, August 1830, located in the HDC.'' The reference is to PDF 84 of the 1830 JH, but the Newel Knight account is cut and pasted from some other source. Clearly the language between Knight's account and the HC are pretty similar, so it is possible that one gave rise to the other. In Vogel's TCHC, he doesn't note that this bit came from Knight or from anywhere else. (3/14/2022: I think I've solved this---I'm guessing the reference is to an 1883 publication of Newel Knight's journal in the \textit{Faith-Promoting Series}. Not totally sure, but see the note in the \hyperref[EXSS-DC27-JSR]{JSR extended commentary on DC 27}, at the part where it talks about Knight. The original manuscript source could be a document called ``Newel Knight autobiography, circa 1871,'' in CHL, \href{https://catalog.churchofjesuschrist.org/record/b99838e6-d9d0-41b5-82f0-66901ac38a8b/0?view=browse}{MS 19156}.)
							}
				\end{quote}
				
			\noindent ``Newel Knight's statement that there were two revelations stands in sharp contrast to the Prophet's that there was only one. Newel Knight further separated these two when he said the first was given in Harmony, Pennsylvania, and the second in Fayette, New York. The difference between these two accounts may be the reason why there is such a variety of dates suggested for this revelation'' (393--394).
			
	% -----
	\subsubsection{Versions}
	% -----
		\label{DC27-V}
		
		See the \href{https://drive.google.com/file/d/1goAvNz8jS4R8slYfMG_db55Ty2z_vmgK/view?usp=sharing}{sections comparison} document.
	
		\begin{compactdesc}
			\item[JSP] \hyperref[EMS-1830-JS-CIR-AUG-1830]{Circa August 1830}; \hyperref[EMS-1835-JS-CIR-AUG-1835]{Circa August 1835}
			\item[BC] Chapter 28, 60
			\item[EMS] 1:10 (March 1833), 78
			\item[DC 1835] Section 50, 179--181
			\item[TS] 4:8 (1 March 1843), 117--118
			\item[MS] 4:10 (February 1844), 151
			\item[DC 1844] Section 50, 270--272
		\end{compactdesc}

	% -----
	\subsubsection{Section 27:1} % *DC27-1 
	% -----
		\label{DC27-1}

		LISTEN to the voice of Jesus Christ, your Lord, your God, and your Redeemer, whose word is quick and powerful.%
			\footnote{%
				% \label{}%
				See \hyperref[HEBREWS-4-12]{Hebrews 4:12}.
				}

		\begin{framed}
		\scriptsize
			\begin{compactdesc}
				\item[JSP] Listen to the voice of Jesus Christ your Lord your God \& your Redeemer whose word is quick \& powerful
				% \item[BC] 
				% \item[]
			\end{compactdesc}
		\end{framed}

	% -----
	\subsubsection{Section 27:2} % *DC27-2 
	% -----
		\label{DC27-2}

		For, behold, I say unto you, that it mattereth not what ye shall eat or what ye shall drink when ye partake of the sacrament, if it so be that ye do it with an eye \hyperref[SINGLE]{single}%
			\footnote{%
				% \label{}%
				The Greek word would be \hyperref[HAPLOUS]{\grl{ἁπλοῦς}}.
				}
		to my glory%
			\footnote{%
				% \label{}%
				This is an important verse from the pragmatic perspective, since Christ is saying that the only physical aspects of the ordinance that actually matter are (1) having something you can eat and drink, and (2) having a certain attitude about the ordinance. Pretty interesting I think.
				}%
		---remembering unto the Father my body which was laid down for you, and my blood which was shed for the remission of your sins.%
			\footnote{%
				% \label{}%
				See \hyperref[MATTHEW-26-26]{Matthew 26:26--29}.
				} 

		\begin{framed}
		\scriptsize
			\begin{compactdesc}
				\item[JSP] for Behold I say unto you it mattereth not what ye <shall> eat or what ye shall drink when ye partake of the sacrament if it so be that ye do it with an eye single to my glory Remembering unto the father my Body which [was] laid down for you \& my blood which was shed for \sout{you} the Remission of your sins
				% \item[BC] 
				% \item[]
			\end{compactdesc}
		\end{framed}

	% -----
	\subsubsection{Section 27:3} % *DC27-3 
	% -----
		\label{DC27-3}

		Wherefore, a commandment I give unto you, that you shall not purchase wine neither strong drink of your enemies;

		\begin{framed}
		\scriptsize
			\begin{compactdesc}
				\item[JSP] Wherefore a commandment I give unto you that ye shall not Purchase Wine neither strong drink of your enemies
				% \item[BC] 
				% \item[]
			\end{compactdesc}
		\end{framed}

	% -----
	\subsubsection{Section 27:4} % *DC27-4 
	% -----
		\label{DC27-4}

		Wherefore, you shall partake of none except it is made new among you;%
			\footnote{%
				% \label{}%
				Here I'm reminded of the ``new wine into old bottles'' teaching of Christ; see \hyperref[MATTHEW-9-17]{Matthew 9:17}, \hyperref[MARK-2-22]{Mark 2:22}, and \hyperref[LUKE-5-37]{Luke 5:37--39}. New wine has to be put in new bottles.
				}
		yea, in this my Father's kingdom which shall be built up on the earth.%
			\footnote{%
				% \label{}%
				This verse is a pretty clear reference to \hyperref[MATTHEW-26-29]{Matthew 26:29}. You could take it as fulfilling the prophecy that Christ makes in that verse, but he seems to say that he himself will be drinking the new wine, and I'm not sure that this verse or some other event in the restoration would really be a fulfillment of that part of the prophecy. Christ does address that part of his prophecy in the next verse, when he promises to drink ``of the fruit of the vine with you on the earth.'' He says that time is coming.
				}

		\begin{framed}
		\scriptsize
			\begin{compactdesc}
				\item[JSP] Wherefore ye shall partake none except it is made new among you yea in this my Fathers Kingdom which shall be built up on the earth
				% \item[BC] 
				% \item[]
			\end{compactdesc}
		\end{framed}

	% -----
	\subsubsection{Section 27:5} % *DC27-5 
	% -----
		\label{DC27-5}

		Behold, this is wisdom in me; wherefore, marvel not, for the hour cometh that I will drink of the fruit of the vine with you on the earth,%
			\footnote{%
				% \label{}%
				See notes in the previous verse.
				}
			\footnote{%
				% \label{}%
				Beginning right here, the 1835 DC begins to add a lot of material that is not in the JSP/RB1 account or the BC account. See the notes above at the beginning of this section and in the \hyperref[DC27-HB]{historical background}.
				}%
		and with Moroni, whom I have sent unto you to reveal the Book of Mormon, containing the fulness of my everlasting gospel, to whom I have committed the keys of the record of the stick of Ephraim;%
			\footnote{%
				% \label{}%
				A reference to \hyperref[EZEKIEL-37-16]{Ezekiel 37:16--19}.
				}

		\begin{framed}
		\scriptsize
			\begin{compactdesc}
				\item[JSP] Behold this is wisdom in me Wherefore marvel not for the hour cometh that I will drink of the fruit of the Vine with you on the Earth
				% \item[BC] 
				% \item[]
			\end{compactdesc}
		\end{framed}

	% -----
	\subsubsection{Section 27:6} % *DC27-6 
	% -----
		\label{DC27-6}

		And also with Elias, to whom I have committed the keys of bringing to pass the restoration of all things spoken by the mouth of all the holy prophets since the world began, concerning the last days;

		\begin{framed}
		\scriptsize
			\begin{compactdesc}
				\item[JSP] ---
				% \item[BC] 
				% \item[]
			\end{compactdesc}
		\end{framed}

	% -----
	\subsubsection{Section 27:7} % *DC27-7 
	% -----
		\label{DC27-7}

		And also John the son of Zacharias, which Zacharias he (Elias)%
			\footnote{%
				% \label{}%
				I'm confused about this---the reference is to \hyperref[LUKE-1-11]{Luke 1:11--20}, but in verse 19 the text is clear that the angel is Gabriel. I think the key here, at least to start thinking about the confusing way that JS uses the term ``Elias,'' is that he uses it to refer to specific individuals---Elijah, for example---but also as a title which numerous individuals take on at different points. Gabriel must be one of those individuals, I think. See the \hyperref[EXSS-DC27-JSR-ELIAS]{extended study in JSR} for more about Elias.
				}
		visited and gave promise that he should have a son, and his name should be John, and he should be filled with the spirit of Elias;

		\begin{framed}
		\scriptsize
			\begin{compactdesc}
				\item[JSP] ---
				% \item[BC] 
				% \item[]
			\end{compactdesc}
		\end{framed}

	% -----
	\subsubsection{Section 27:8} % *DC27-8 
	% -----
		\label{DC27-8}

		Which John I have sent unto you, my servants, Joseph Smith, Jun., and Oliver Cowdery, to ordain you unto the first priesthood which you have received, that you might be called and ordained even as Aaron;

		\begin{framed}
		\scriptsize
			\begin{compactdesc}
				\item[JSP] ---
				% \item[BC] 
				% \item[]
			\end{compactdesc}
		\end{framed}

	% -----
	\subsubsection{Section 27:9} % *DC27-9 
	% -----
		\label{DC27-9}

		And also Elijah, unto whom I have committed the keys of the power of turning the hearts of the fathers to the children, and the hearts of the children to the fathers, that the whole earth may not be smitten with a curse;

		\begin{framed}
		\scriptsize
			\begin{compactdesc}
				\item[JSP] ---
				% \item[BC] 
				% \item[]
			\end{compactdesc}
		\end{framed}

	% -----
	\subsubsection{Section 27:10} % *DC27-10 
	% -----
		\label{DC27-10}

		And also with Joseph and Jacob, and Isaac, and Abraham, your fathers, by whom the promises remain;

		\begin{framed}
		\scriptsize
			\begin{compactdesc}
				\item[JSP] ---
				% \item[BC] 
				% \item[]
			\end{compactdesc}
		\end{framed}

	% -----
	\subsubsection{Section 27:11} % *DC27-11 
	% -----
		\label{DC27-11}

		And also with Michael, or Adam,%
			\footnote{%
				% \label{}%
				See discussion of Michael/Adam at the \hyperref[EXSS-DC27-JSR]{extended study from JSR} for this section.
				}
		the father of all, the prince of all, the ancient of days;

		\begin{framed}
		\scriptsize
			\begin{compactdesc}
				\item[JSP] ---
				% \item[BC] 
				% \item[]
			\end{compactdesc}
		\end{framed}

	% -----
	\subsubsection{Section 27:12} % *DC27-12 
	% -----
		\label{DC27-12}

		And also with Peter, and James, and John, whom I have sent unto you, by whom I have ordained you and confirmed you to be apostles,%
			\footnote{%
				% \label{}%
				This is the record we have of the restoration of the Melchizedek priesthood. There definitely does seem to be something in common with temple ordinances, or the endowment ceremony, here.
				}
		and especial witnesses of my name, and bear the keys of your ministry and of the same things which I revealed unto them;

		\begin{framed}
		\scriptsize
			\begin{compactdesc}
				\item[JSP] ---
				% \item[BC] 
				% \item[]
			\end{compactdesc}
		\end{framed}

	% -----
	\subsubsection{Section 27:13} % *DC27-13 
	% -----
		\label{DC27-13}

		Unto whom I have committed the keys of my kingdom, and a dispensation of the gospel for the last times; and for the fulness of times, in the which I will gather together in one all things, both which are in heaven, and which are on earth;

		\begin{framed}
		\scriptsize
			\begin{compactdesc}
				\item[JSP] ---
				% \item[BC] 
				% \item[]
			\end{compactdesc}
		\end{framed}

	% -----
	\subsubsection{Section 27:14} % *DC27-14 
	% -----
		\label{DC27-14}

		And also with all those whom my Father hath given me out of the world.

		\begin{framed}
		\scriptsize
			\begin{compactdesc}
				\item[JSP] \& with \sout{you} all those whom my father hath given me out of the world
				% \item[BC] 
				% \item[]
			\end{compactdesc}
		\end{framed}

	% -----
	\subsubsection{Section 27:15} % *DC27-15 
	% -----
		\label{DC27-15}

		Wherefore, lift up your hearts and rejoice, and gird up your loins, and take upon you my whole armor,%
			\footnote{%
				% \label{}%
				For this verse and below, see \hyperref[EPHESIANS-6-11]{Ephesians 6:11--18}.
				}
		that ye may be able to withstand the evil day, having done all, that ye may be able to stand.

		\begin{framed}
		\scriptsize
			\begin{compactdesc}
				\item[JSP] Wherefore lift up your hearts \& rejoice \& Gird up your loins
				% \item[BC] 
				% \item[]
			\end{compactdesc}
		\end{framed}

	% -----
	\subsubsection{Section 27:16} % *DC27-16 
	% -----
		\label{DC27-16}

		Stand, therefore, having your loins girt about with truth, having on the breastplate of righteousness, and your feet shod with the preparation of the gospel of peace, which I have sent mine angels to commit unto you;

		\begin{framed}
		\scriptsize
			\begin{compactdesc}
				\item[JSP] ---
				% \item[BC] 
				% \item[]
			\end{compactdesc}
		\end{framed}

	% -----
	\subsubsection{Section 27:17} % *DC27-17 
	% -----
		\label{DC27-17}

		Taking the shield of faith wherewith ye shall be able to quench all the fiery darts of the wicked;

		\begin{framed}
		\scriptsize
			\begin{compactdesc}
				\item[JSP] ---
				% \item[BC] 
				% \item[]
			\end{compactdesc}
		\end{framed}

	% -----
	\subsubsection{Section 27:18} % *DC27-18 
	% -----
		\label{DC27-18}

		And take the helmet of salvation, and the sword of my Spirit, which I will pour out upon you, and my word which I reveal unto you, and be agreed as touching all things whatsoever ye ask of me, and be faithful until I come, and ye shall be caught up, that where I am ye shall be also.  Amen.

		\begin{framed}
		\scriptsize
			\begin{compactdesc}
				\item[JSP] \& be \sout{faitful} faithful untill I come even so amen
				% \item[BC] 
				% \item[]
			\end{compactdesc}
		\end{framed}